% =====================================================================
% Appendix B — Hyperparameters + training trajectory
% =====================================================================

\section{Hyperparameters and training trajectory}
\label{app:hparams}

\begin{table}[h]
\centering\small
\caption{Baseline training hyperparameters.}
\label{tab:hparams}
\begin{tabular}{@{}ll@{}}
\toprule
History length $S$            & 24 hours                                    \\
Future horizon $H$            & 24 hours                                    \\
Spatial token grid            & $8\times 8$ ($P=64$)                        \\
Embedding dim $D$             & 128                                         \\
Encoder layers                & 4                                           \\
Attention heads               & 4                                           \\
Dropout                       & 0.1                                         \\
Optimizer                     & AdamW \citep{loshchilov2019adamw}            \\
Initial LR                    & $3\times 10^{-4}$                           \\
Weight decay                  & $10^{-2}$                                   \\
Schedule                      & cosine warm-restarts \citep{loshchilov2017sgdr} \\
Batch size                    & 16                                          \\
Epochs                        & 14                                          \\
Loss                          & MSE in z-score space                        \\
Mixed precision               & PyTorch AMP fp16                            \\
Random seeds                  & not set (caveat below)                      \\
\bottomrule
\end{tabular}
\end{table}

\paragraph{Random seeds caveat.}
The training pipeline does not call \texttt{torch.manual\_seed} or
\texttt{np.random.seed}. The $5.24~\%$ test MAPE number reported in this
paper is pinned to the specific checkpoint
\texttt{space/checkpoints/best.pt} (and equivalently
\texttt{runs/cnn\_transformer\_baseline/checkpoints/best.pt}), which
we ship with the public release; it is not bit-reproducible across
re-training runs. The deployed numbers in §\ref{sec:drift} and the
multi-window experiment are reproducible (same checkpoint, same code,
same data).

\paragraph{Per-epoch validation MAPE (val 2021).}

\begin{table}[h]
\centering\small
\caption{Per-epoch validation MAPE on 2021 data; baseline only.}
\label{tab:trajectory}
\begin{tabular}{@{}rr@{}}
\toprule
Epoch & val MAPE \\
\midrule
0 & 11.22 \% \\
4 & 9.16 \%  \\
6 & 8.76 \%  \\
13 (final) & \textbf{6.92 \%} \\
\bottomrule
\end{tabular}
\end{table}

The training procedure ran continuously for 14 epochs (no chained-resume
LR-scheduler reset), reaching the small-LR fine-tuning regime where the
final validation drop occurs.
